\studySubsection{calc-02-sequence-limit}{第 2 讲 数列极限}

\begin{knowledgeBox}
教材页码：第 76 页起。

本讲用于整理数列极限、单调有界准则、夹逼准则、递推数列极限和常见数列变形。新增题目时重点检查极限存在性、递推关系稳定点和放缩方向。
\end{knowledgeBox}

% BEGIN calculus-foundation:calc-02-pages
\knowledgeAnchor[MATH1-KN-CALC-0065]{数列极限的 $\varepsilon$-$N$ 定义}
\begin{knowledgeBox}
\textbf{定义或结论：} $a_n\to A$ 指对任意 $\varepsilon>0$，存在 $N$，当 $n>N$ 时有 $|a_n-A|<\varepsilon$。

\begin{itemize}
  \item \textbf{直觉与必会动作：} 能按定义解释“任意精度最终都进入邻域”；会做最基础的定义证明。
  \item \textbf{成立条件：} $N$ 可依赖于 $\varepsilon$，但不能依赖于具体的 $n$；“任意”与“存在”的顺序不可交换。
  \item \textbf{典型用法与识别信号：} 按定义证明简单数列极限；否定极限定义；识别信号为 “按定义”“对任意ε”“存在N”。
  \item \textbf{关系与迁移：} 前置知识为 绝对值不等式；可迁移到 收敛性质、反例构造、选择题逻辑。
  \item \textbf{常见误区：} 量词顺序错误；把“存在ε”写成“任意ε”。迁移到新题时先复核定义域、趋近方向和所用定理的条件。
\end{itemize}
\end{knowledgeBox}

\knowledgeAnchor[MATH1-KN-CALC-0066]{收敛数列的唯一性}
\begin{knowledgeBox}
\textbf{定义或结论：} 若 $a_n\to A$ 且 $a_n\to B$，则 $A=B$。

\begin{itemize}
  \item \textbf{直觉与必会动作：} 理解一个数列不可能收敛到两个不同极限。
  \item \textbf{成立条件：} 前提是同一数列在同一指标方向 $n\to\infty$。
  \item \textbf{典型用法与识别信号：} 判断参数；证明极限唯一；识别信号为 “极限存在且等于”。
  \item \textbf{关系与迁移：} 前置知识为 极限定义；可迁移到 递推数列极限方程的候选筛选。
  \item \textbf{常见误区：} 由递推方程得到多个根后，误认为都是极限。迁移到新题时先复核定义域、趋近方向和所用定理的条件。
\end{itemize}
\end{knowledgeBox}

\knowledgeAnchor[MATH1-KN-CALC-0067]{收敛数列有界性与保号性}
\begin{knowledgeBox}
\textbf{定义或结论：} $a_n\to A$ 则 $(a_n)$ 有界；若 $A>0$，则充分大 $n$ 有 $a_n>0$。

\begin{itemize}
  \item \textbf{直觉与必会动作：} 掌握收敛必有界；极限为正时充分靠后项为正。
  \item \textbf{成立条件：} 有界是收敛的必要非充分条件；保号性只保证“最终”，不能保证所有项。
  \item \textbf{典型用法与识别信号：} 判断命题真假；估计分母不为零；识别信号为 “收敛必…”“充分大n”。
  \item \textbf{关系与迁移：} 前置知识为 极限定义；可迁移到 极限运算、级数必要条件、递推数列。
  \item \textbf{常见误区：} 把有界当收敛；把最终性质推广到前有限项。迁移到新题时先复核定义域、趋近方向和所用定理的条件。
\end{itemize}
\end{knowledgeBox}

\knowledgeAnchor[MATH1-KN-CALC-0068]{子列与聚点判据}
\begin{knowledgeBox}
\textbf{定义或结论：} 若 $a_n\to A$，则任意子列 $a_{n_k}\to A$；若找到两子列极限不同，则原数列发散。

\begin{itemize}
  \item \textbf{直觉与必会动作：} 理解收敛数列任一子列同极限；会用两子列不同极限证明发散。
  \item \textbf{成立条件：} 两子列必须指标严格递增且趋于无穷；存在一个收敛子列不能推出原数列收敛。
  \item \textbf{典型用法与识别信号：} 振荡数列判敛；构造反例；识别信号为 “偶数项/奇数项”“振荡”。
  \item \textbf{关系与迁移：} 前置知识为 数列极限；可迁移到 函数极限的序列判据、级数部分和。
  \item \textbf{常见误区：} 只看一个子列；把部分项趋势当整体趋势。迁移到新题时先复核定义域、趋近方向和所用定理的条件。
\end{itemize}
\end{knowledgeBox}

\knowledgeAnchor[MATH1-KN-CALC-0069]{数列极限四则运算}
\noindent{\footnotesize\textbf{索引：} \knowledgeIndexRef[MATH1-KN-CALC-0069]{相邻幂差的降阶等价}}\par
\begin{knowledgeBox}
\textbf{定义或结论：} 若 $a_n\to A,b_n\to B$，则和、差、积相应收敛；若 $B\ne0$，则 $a_n/b_n\to A/B$。

\begin{itemize}
  \item \textbf{直觉与必会动作：} 熟练处理和差积商与常数倍；商法则先检查分母极限非零。
  \item \textbf{成立条件：} 不能对 $0/0$、$\infty/\infty$ 直接分别取极限；无穷大不是普通实数。
  \item \textbf{典型用法与识别信号：} 有理式数列极限；拆分极限；识别信号为 “各项极限已知”。
  \item \textbf{关系与迁移：} 前置知识为 极限性质；可迁移到 夹逼、等价、递推数列。
  \item \textbf{常见误区：} 在分母极限为0时套商法则；把不存在的分项极限拆开。迁移到新题时先复核定义域、趋近方向和所用定理的条件。
\end{itemize}

\textbf{相邻幂差的降阶等价：} 对固定实数 $k\ne0$，当 $n\to\infty$ 时，
\[
n^k-(n-1)^k\sim k n^{k-1}.
\]
相邻差可看成步长为 $1$ 的离散求导。严格地，由拉格朗日中值定理，存在 $\xi_n\in(n-1,n)$，使
\[
n^k-(n-1)^k=k\xi_n^{k-1}\sim k n^{k-1}.
\]
当 $k$ 为正整数时，也可由二项式定理
\[
(n-1)^k=n^k-kn^{k-1}+\binom{k}{2}n^{k-2}-\cdots
\]
直接看出：两个 $n^k$ 项抵消后，第一个非零项为 $k n^{k-1}$。更一般地，对固定 $h$ 有
\[
n^k-(n-h)^k\sim kh n^{k-1}.
\]
\end{knowledgeBox}

\problemAnchor{MATH1-CALC-0012}
\begin{problemBox}
\begin{problemMeta}
\textbf{题目编号：} \problemRef{MATH1-CALC-0012} \quad
\textbf{索引：} \problemIndexRef{MATH1-CALC-0012} \quad
\textbf{来源：} 用户图片 / 教材例 2.11（书名未注明）

\textbf{题型：} 隐式约束下耦合数列的极限与等价关系 \quad
\textbf{难度：} 中等

\textbf{知识点：} \knowledgeRef[MATH1-KN-CALC-0070]{夹逼准则}，
\knowledgeRef[MATH1-KN-CALC-0081]{常用等价无穷小}，
\knowledgeIndexRef[MATH1-KN-CALC-0006]{方程组消元}
\end{problemMeta}

\textbf{题目：}

设
\[
0<a_n<\frac{\pi}{2},\qquad 0<b_n<\frac{\pi}{2},\qquad
\cos a_n-a_n=\cos b_n,
\]
且 \(b_n\to0\)。求
\[
\lim_{n\to\infty}a_n,\qquad
\lim_{n\to\infty}\frac{a_n}{b_n^2}.
\]
\end{problemBox}

\begin{solutionBox}
\textbf{共同的第一步：先确定趋向。}

由题设
\[
\cos a_n-\cos b_n=a_n>0.
\]
函数 \(\cos x\) 在 \(\left(0,\frac{\pi}{2}\right)\) 上严格递减，故
\[
0<a_n<b_n.
\]
再由 \(b_n\to0\) 及夹逼准则，得
\[
\boxed{a_n\to0}.
\]

\textbf{方法一：选取桥梁量，按约束渐近消元。}

目标分母为 \(b_n^2\)，自然联想到
\(1-\cos b_n\sim\frac12b_n^2\)。原关系恰可改写为
\[
1-\cos b_n=(1-\cos a_n)+a_n.
\]
于是
\begin{align*}
\frac{a_n}{b_n^2}
&=\frac{1-\cos b_n}{b_n^2}
  \cdot\frac{a_n}{1-\cos b_n}\\
&=\frac{1-\cos b_n}{b_n^2}
  \cdot\frac{1}{\dfrac{1-\cos a_n}{a_n}+1}.
\end{align*}
其中
\[
\frac{1-\cos b_n}{b_n^2}\to\frac12,
\qquad
\frac{1-\cos a_n}{a_n}
=\frac{1-\cos a_n}{a_n^2}\,a_n\to\frac12\cdot0=0.
\]
故
\[
\boxed{\lim_{n\to\infty}\frac{a_n}{b_n^2}=\frac12}.
\]

这里并没有求出精确的 \(a_n=g(b_n)\)，而是只利用约束消去了求极限所需的主导部分。

\textbf{方法二：和差化积，同时得到阶的估计。}

令
\[
U_n=\frac{a_n+b_n}{2},\qquad V_n=\frac{b_n-a_n}{2}.
\]
则 \(U_n,V_n\to0\)，且
\[
a_n=\cos a_n-\cos b_n=2\sin U_n\sin V_n.
\]
由 \(0<\sin x\le x\) 得
\[
0<a_n\le 2U_nV_n
=\frac{b_n^2-a_n^2}{2}\le\frac{b_n^2}{2},
\]
从而 \(a_n/b_n\to0\)。再写成
\[
\frac{a_n}{b_n^2}
=\frac12\cdot\frac{\sin U_n}{U_n}
\cdot\frac{\sin V_n}{V_n}
\cdot\left(1-\frac{a_n^2}{b_n^2}\right),
\]
各因子取极限即得 \(1/2\)。

\textbf{方法三：先估阶，再作泰勒展开。}

由
\[
a_n=(1-\cos b_n)-(1-\cos a_n)
\]
可知
\[
0<a_n\le1-\cos b_n\le\frac{b_n^2}{2},
\]
所以 \(a_n=O(b_n^2)\)，进而 \(a_n^2=o(b_n^2)\)。于是
\begin{align*}
a_n
&=\cos a_n-\cos b_n\\
&=\frac{b_n^2-a_n^2}{2}+o(b_n^2)
=\frac12b_n^2+o(b_n^2),
\end{align*}
故仍有 \(a_n/b_n^2\to1/2\)。

\textbf{答案：}
\[
\boxed{\lim_{n\to\infty}a_n=0,\qquad
\lim_{n\to\infty}\frac{a_n}{b_n^2}=\frac12}.
\]
\end{solutionBox}

\begin{knowledgeBox}
\textbf{通用思想：利用约束降维，并在极限层面消元。}

若若干变量或数列由关系 \(F(u_n,v_n)=0\) 耦合，不必总把其中一个精确解成另一个的函数；极限题往往只需找出足以确定目标的渐近关系。标准流程是：
\begin{enumerate}
  \item \textbf{先定趋向：} 用定义域、单调性、保号性或夹逼确定各量趋向何处；
  \item \textbf{再定阶：} 由约束建立 \(O\)、\(o\) 或等价关系，判断谁是主导项；
  \item \textbf{最后定系数：} 选取与目标同形的桥梁量，把目标拆成标准极限与约束因子的乘积。
\end{enumerate}

这比“强行化为一个变量”更一般：它既包含精确代换、方程组消元，也包含只在主导阶上成立的渐近消元。若约束不足，则不能如此降维；例如仅知 \(u_n+v_n\to0\)，并不能唯一确定 \(u_n/v_n\) 的极限。使用前必须确认约束与附加条件足以排除不同路径。

可记为：\textbf{先定趋向，再定阶，最后定系数。}
\end{knowledgeBox}

\problemAnchor{MATH1-CALC-0013}
\begin{problemBox}
\begin{problemMeta}
\textbf{题目编号：} \problemRef{MATH1-CALC-0013} \quad
\textbf{索引：} \problemIndexRef{MATH1-CALC-0013} \quad
\textbf{来源：} 用户图片 / 题号 2.4（书名未注明）

\textbf{题型：} 相邻幂差判阶与参数极限 \quad
\textbf{难度：} 基础

\textbf{知识点：} \knowledgeRef[MATH1-KN-CALC-0069]{数列极限四则运算与相邻幂差降阶}
\end{problemMeta}

\textbf{题目：}

设 $k$ 为实常数。若有限极限
\[
\lim_{n\to\infty}\frac{n^{99}}{n^k-(n-1)^k}
\]
存在且不为零，求 $k$。
\end{problemBox}

\begin{solutionBox}
\textbf{思路入口：先判断分母的真实阶数。}

分母虽然由两个 $k$ 次量组成，但它们的最高阶项会相消。$k=0$ 时分母恒为零，应先排除。对 $k\ne0$，由相邻幂差的降阶等价，
\[
n^k-(n-1)^k\sim k n^{k-1}.
\]
因此
\[
\frac{n^{99}}{n^k-(n-1)^k}
\sim \frac{n^{99}}{k n^{k-1}}
=\frac1k n^{100-k}.
\]
要使极限为有限非零常数，必须有
\[
100-k=0,
\]
故
\[
\boxed{k=100}.
\]
此时原极限为
\[
\boxed{\frac1{100}}.
\]

若题目背景把 $k$ 限定为正整数，也可以直接使用二项式定理：
\[
n^k-(n-1)^k
=k n^{k-1}-\binom{k}{2}n^{k-2}+\cdots.
\]
分母的有效最高次数是 $k-1$，与分子的 $99$ 次匹配便得到 $k-1=99$。
\end{solutionBox}

\begin{mistakeBox}
\textbf{本题实际卡点：} 直觉上把分母看成 $k$ 次量，于是猜测 $k=99$。

\textbf{第一个失效步骤：} 未先展开或估计 $n^k-(n-1)^k$，忽略了两个 $n^k$ 主项完全抵消；分母的真实阶数是 $k-1$，不是 $k$。

\textbf{下次检查动作：} 看到两个同阶大项相减，先问“最高阶项是否抵消”，再匹配次数。可用二项式展开、中值定理或泰勒展开找到相消后的第一个非零项。
\end{mistakeBox}

\problemAnchor{MATH1-CALC-0014}
\begin{problemBox}
\begin{problemMeta}
\textbf{题目编号：} \problemRef{MATH1-CALC-0014} \quad
\textbf{索引：} \problemIndexRef{MATH1-CALC-0014} \quad
\textbf{来源：} 用户图片 / 题号 2.7（书名未注明）

\textbf{题型：} 闭区间有界性与 $n$ 次方根夹逼 \quad
\textbf{难度：} 基础

\textbf{知识点：} \knowledgeRef[MATH1-KN-CALC-0091]{有界性与最大最小值定理}，
\knowledgeRef[MATH1-KN-CALC-0070]{夹逼准则}
\end{problemMeta}

\textbf{题目：}

设函数 $f(x)$ 在 $[a,b]$ 上连续，点列
$\{x_k\}_{k\ge1}\subset[a,b]$。求
\[
\lim_{n\to\infty}
\sqrt[n]{\frac1n\sum_{k=1}^{n}e^{f(x_k)}}.
\]
\end{problemBox}

\begin{solutionBox}
\textbf{思路入口：不求平均值本身的极限，只寻找与 $n$ 无关的正上、下界。}

这里使用闭区间连续函数的最大最小值定理：若实值函数在有限闭区间上连续，
则它不仅有上、下界，而且能在区间内取得有限的最大值和最小值。闭区间与连续性
共同保证这个结论；若去掉其中一个条件，一般不能保证有界。于是可记
\[
m=\min_{x\in[a,b]}f(x),\qquad
M=\max_{x\in[a,b]}f(x).
\]
对每个 $x_k\in[a,b]$，由指数函数严格递增可得
\[
e^m\le e^{f(x_k)}\le e^M.
\]
逐项求和再除以 $n$，得到
\[
e^m\le
\frac1n\sum_{k=1}^{n}e^{f(x_k)}
\le e^M.
\]
三项均为正数，开 $n$ 次方后不等号方向不变：
\[
e^{m/n}\le
\sqrt[n]{\frac1n\sum_{k=1}^{n}e^{f(x_k)}}
\le e^{M/n}.
\]
因为 $m,M$ 是与 $n$ 无关的有限常数，故
\[
\frac{m}{n}\to0,\qquad \frac{M}{n}\to0.
\]
指数函数在 $0$ 处连续，所以
\[
e^{m/n}\to e^0=1,\qquad e^{M/n}\to e^0=1.
\]
由夹逼准则，原极限为
\[
\boxed{1}.
\]

点列 $\{x_k\}$ 是否收敛、如何分布都不重要；闭区间上的统一界已经足以决定这个极限。
\end{solutionBox}

\begin{knowledgeBox}
\textbf{可复用结论：} 若存在与 $n$ 无关的常数 $0<c\le C$，使
\[
c\le u_n\le C,
\]
则
\[
u_n^{1/n}\to1.
\]
本题中取
\[
u_n=\frac1n\sum_{k=1}^{n}e^{f(x_k)},
\qquad c=e^m,\quad C=e^M.
\]

必须注意“正下界”不能省略。仅知道 $u_n>0$ 且有界并不足以推出结论，例如
$u_n=e^{-n}$ 有界，但 $u_n^{1/n}=e^{-1}$。更本质的判据是
$\ln u_n=o(n)$。
\end{knowledgeBox}

\studySubsection{calc-02-convergence-criteria}{数列极限判据与递推数列}

\begin{knowledgeBox}
本页把“候选极限”和“极限存在性”严格分开：不动点方程只负责找候选，夹逼或单调有界负责证明收敛。
\end{knowledgeBox}

\knowledgeAnchor[MATH1-KN-CALC-0070]{夹逼准则}
\begin{knowledgeBox}
\textbf{定义或结论：} 若从某项起 $x_n\le y_n\le z_n$ 且 $x_n,z_n\to A$，则 $y_n\to A$。

\begin{itemize}
  \item \textbf{直觉与必会动作：} 会构造上下界并证明两端同极限。
  \item \textbf{成立条件：} 不等式只需从某项起成立；两端必须趋于同一有限或同向无穷极限。
  \item \textbf{典型用法与识别信号：} 含和式、积分式、振荡因子的数列极限；识别信号为 “夹在”“绝对值≤”“和式难算”。
  \item \textbf{关系与迁移：} 前置知识为 不等式、极限运算；可迁移到 定积分估计、级数余项、概率极限直觉。
  \item \textbf{常见误区：} 上、下界极限不同；不说明不等式成立范围。迁移到新题时先复核定义域、趋近方向和所用定理的条件。
\end{itemize}
\end{knowledgeBox}

\knowledgeAnchor[MATH1-KN-CALC-0071]{单调有界准则}
\begin{knowledgeBox}
\textbf{定义或结论：} 单调递增且上有界，或单调递减且下有界的数列必收敛。

\begin{itemize}
  \item \textbf{直觉与必会动作：} 会证明递推数列单调且有界，从而得到收敛性。
  \item \textbf{成立条件：} 必须同时证明单调与相应方向的界；仅由递推极限方程不能证明极限存在。
  \item \textbf{典型用法与识别信号：} 递推数列极限；证明数列收敛；识别信号为 “递推定义”“证明极限存在”。
  \item \textbf{关系与迁移：} 前置知识为 数学归纳法、单调性；可迁移到 不等式、函数单调性、积分递推。
  \item \textbf{常见误区：} 先设极限再代入，跳过存在性；界选得不封闭。迁移到新题时先复核定义域、趋近方向和所用定理的条件。
\end{itemize}
\end{knowledgeBox}

\knowledgeAnchor[MATH1-KN-CALC-0072]{递推数列极限}
\begin{knowledgeBox}
\textbf{定义或结论：} 若已知 $a_n\to A$ 且递推函数连续，则 $A=\varphi(A)$；此方程只给候选，不给存在性。

\begin{itemize}
  \item \textbf{直觉与必会动作：} 会用不动点方程求候选，并用单调有界、压缩或夹逼证明存在。
  \item \textbf{成立条件：} 取极限前须证明收敛或题目已给收敛；多个不动点需结合初值与不变区间筛选。
  \item \textbf{典型用法与识别信号：} 求递推极限；证明单调有界；识别信号为 “$a_{n+1}=\varphi(a_n)$”。
  \item \textbf{关系与迁移：} 前置知识为 单调有界、连续性；可迁移到 中值定理、积分不等式、函数迭代。
  \item \textbf{常见误区：} 循环论证；忽略初值导致选错根。迁移到新题时先复核定义域、趋近方向和所用定理的条件。
\end{itemize}
\end{knowledgeBox}

\knowledgeAnchor[MATH1-KN-CALC-0314]{递推数列极限}
\begin{knowledgeBox}
\textbf{题型识别：} 给初值和 $a_{n+1}=\varphi(a_n)$，求极限或证收敛；识别信号为 “递推定义”“证明数列收敛”“求lim a\_n”。

\begin{itemize}
  \item \textbf{标准流程：} ①解不动点方程找候选；②由初值构造不变区间；③证明单调或压缩；④由单调有界得收敛；⑤代入连续递推求极限；⑥排除其他根。
  \item \textbf{方法选择：} 能证单调有界时最规范；有收缩比时用误差估计；可替代路线包括 夹逼、误差递推、数学归纳。
  \item \textbf{合法性检查：} 每一步先核对定义域、极限或定理条件；特别防止 先设极限再代入，循环论证；不筛多个不动点。
  \item \textbf{考场步骤：} 先写识别出的结构与必要条件，再按标准流程逐步化简，最后回代、筛根或复核单侧情形。
  \item \textbf{迁移方向：} 分式递推、根式递推、积分递推；可与 极限、单调性、不等式 交错训练。
\end{itemize}
\end{knowledgeBox}
% END calculus-foundation:calc-02-pages
